\begingroup
\renewcommand{\figura}[1]{\begin{center}\begin{adjustbox}{max width=\linewidth,max totalheight=4.0cm}#1\end{adjustbox}\end{center}}
\section{Panorama}
\begin{frame}{Microeletrônica: integrar para transformar}
\figura{\begin{tikzpicture}[x=1cm,y=1cm,font=\small]
\draw[blue_theme,thick,fill=white] (3.6,0) rectangle (6.4,2.8);
\foreach \x/\y/\t in {3.8/.3/Memória,5.1/.3/Lógica,3.8/1.6/Analógico,5.1/1.6/RF}{\fill[blue_theme!12] (\x,\y) rectangle ++(1.1,.9);\node[font=\scriptsize] at (\x+.55,\y+.45){\t};}
\node[align=center] (a) at (.7,3.6){Saúde\\\scriptsize sensores e diagnóstico};
\node[align=center] (b) at (9.3,3.6){Comunicação\\\scriptsize rádio e redes};
\node[align=center] (c) at (.7,-.5){Mobilidade\\\scriptsize controle e potência};
\node[align=center] (d) at (9.3,-.5){Computação\\\scriptsize dados e inteligência};
\draw[seta] (3.6,2.5)--(a);\draw[seta] (6.4,2.5)--(b);\draw[seta] (3.6,.3)--(c);\draw[seta] (6.4,.3)--(d);
\path[use as bounding box] (-1,-1.2) rectangle (11,4.3);
\end{tikzpicture}}
\chamada{Projetar dispositivos, circuitos e sistemas em um chip.}
\end{frame}

\begin{frame}{Do componente discreto ao sistema integrado}
\figura{\linhaHistorica}
\chamada{Integrar reduz conexões externas e amplia o que um circuito faz.}
\fonte{\href{https://www.computerhistory.org/siliconengine/timeline/}{Computer History Museum — The Silicon Engine}; \href{https://www.intel.com/content/www/us/en/history/virtual-vault/articles/the-intel-4004.html}{Intel — 4004}. Linha do tempo seletiva; espaçamento não cronológico.}
\end{frame}

% As micrografias só são incluídas quando os arquivos científicos estão locais.
% Não usar ilustrações sintéticas ou fotografias de encapsulamento como substitutas.
\IfFileExists{assets/microscopia/vlsi_nise.jpg}{%
\begin{frame}{Dentro do chip: muitas conexões, muitos dispositivos}
\begin{center}\includegraphics[width=.92\linewidth,height=4.5cm,keepaspectratio]{assets/microscopia/vlsi_nise.jpg}\end{center}
\chamada{A organização do circuito aparece ao ampliar suas camadas.}
\fonte{NISE Network / Lawrence Hall of Science, \href{https://www.nisenet.org/catalog/media/zoom_microchip_video}{Zoom into a Microchip}, recurso de 2013. \href{https://creativecommons.org/licenses/by-nc-sa/3.0/us/}{CC BY-NC-SA 3.0 US}. Imagem estática; circuito mais antigo e processo não identificado.}
\end{frame}}{}
\IfFileExists{assets/microscopia/gaa_imec.jpg}{%
\begin{frame}{Do circuito ao transistor nanométrico: uma vista TEM}
\begin{center}\includegraphics[width=.92\linewidth,height=4.5cm,keepaspectratio]{assets/microscopia/gaa_imec.jpg}\end{center}
\chamada{Nanosheets: canais empilhados envolvidos pelo Gate.}
\fonte{imec, \href{https://www.imec-int.com/en/articles/outer-wall-forksheet-bridge-nanosheet-and-cfet-device-architectures-logic-technology}{Outer wall forksheet: bridging nanosheet and CFET}, Fig. 1 (2025). Microscopia eletrônica de transmissão (TEM). Amostra diferente da vista NISE; não é um zoom contínuo do mesmo chip.}
\end{frame}}{}



\begin{frame}{Lei de Moore: uma trajetória de integração}
\figura{\graficoMoore}
\small\centering 1965: previsão de duplicação anual.\quad 1975: revisão para dois anos.
\chamada{Uma observação histórica que orientou a indústria.}
\fonte{\href{https://www.intel.com/content/www/us/en/newsroom/resources/moores-law.htm}{Intel — Moore's Law}. Curva ideal normalizada, sem dados medidos. Refere-se à quantidade de transistores por circuito integrado, não à frequência nem a uma lei física.}
\end{frame}

% Gerado por scripts/gerar_integracao.py.
\begin{frame}{Integração: de milhares a bilhões de transistores}
\only<1|handout:0>{\figura{\graficoIntegracao{1}}\etapa{1/2}{Marcos Intel: mais transistores por circuito integrado.}}
\only<2|handout:1>{\figura{\graficoIntegracao{2}}\etapa{2/2}{SoCs Apple e GPU NVIDIA ampliam o panorama; funções e áreas são distintas.}}
\fonte{Fontes primárias: Intel (guia histórico e Nehalem), ISSCC 2012 (Ivy Bridge), Apple Newsroom (M1--M4) e NVIDIA (GA100). Eixo vertical logarítmico. Contagens por die, não por placa; processo é a designação publicada, não comprimento de porta. Dados e links em dados/integracao\_processadores.csv.}
\end{frame}
\begin{frame}{Ler o gráfico: o que significa ``integrar mais?}
\duas{\textcolor{blue_theme}{\large\textbf{O que contamos}}\par\medskip Mais dispositivos no mesmo circuito integrado.\par\medskip CPU, memórias, interfaces e aceleradores podem dividir o die.}{\textcolor{blue_theme}{\large\textbf{O que também importa}}\par\medskip Processo, área do die e função do chip.\par\medskip A contagem sozinha não mede desempenho nem densidade.}
\chamada{VLSI: integração de muitos dispositivos e conexões em um chip.}
\fonte{M1--M4 são SoCs com arquitetura Arm; A100 é uma GPU. O histórico é uma seleção didática de produtos, sem ajuste de uma curva universal de Moore. As contagens recentes abrangem o chip completo.}
\end{frame}



\endgroup
